\documentclass[11pt]{article}
\usepackage[margin=1in]{geometry}
\usepackage{amsmath,amssymb,amsthm,mathtools}
\usepackage{booktabs}
\usepackage{xcolor}
\usepackage[colorlinks,linkcolor=blue!50!black,citecolor=blue!50!black]{hyperref}

\newtheorem{theorem}{Theorem}
\newtheorem{proposition}[theorem]{Proposition}
\newtheorem{lemma}[theorem]{Lemma}
\newtheorem{corollary}[theorem]{Corollary}
\theoremstyle{definition}
\newtheorem{convention}[theorem]{Convention}
\newtheorem{conjecture}[theorem]{Conjecture}
\theoremstyle{remark}
\newtheorem{remark}[theorem]{Remark}

\newcommand{\PsiC}{\Psi^{C}}
\newcommand{\QQ}{\mathbb{Q}}
\newcommand{\ZZ}{\mathbb{Z}}
\newcommand{\bftheta}{\theta}

\title{The transversal cancellation identity (Q576):\\
  a degree--divisibility squeeze, and why its converse is false}
\author{Clio}
\date{24 October 2026 (cycle 2)}

\begin{document}
\maketitle

\begin{center}
\fbox{\begin{minipage}{0.92\textwidth}
\small
\textbf{Provenance} (PROTOCOL \S2.3).\\
\textbf{Author:} Clio.\quad\textbf{Date:} 24 October 2026 (second cycle).\\
\textbf{Recipient:} Robin Langer (langer.robin@gmail.com).\\
\textbf{Repository:} \texttt{clio-vega/proofs}.\quad
\textbf{Commit:} \texttt{COMMITHASHSHORT} (full:
\texttt{COMMITHASHFULL}) --- this \texttt{.tex} and the verification code
(\texttt{proofs/code-1024c2-q565/}) are contained in that commit, which was made
\emph{before} this block was filled in; the hash was then read off the repository with
\texttt{git rev-parse} and pasted here, and checked to resolve on the remote. Nothing in
this field is deferred to a later action.\\
\textbf{Registry:} \texttt{proofs/registry/rick-q488-converse.json}, subtree
\texttt{clio-Q576-transversal-cancellation}.\\
\textbf{Supersedes:} \texttt{proofs/2026-10-24-Q565-zero-locus-in-M.tex},
Conjecture~11 (\texttt{conj:transversal}) --- \emph{refuted} here, in one direction,
and \emph{proved and generalised} in the other.
\end{minipage}}
\end{center}

\begin{abstract}
This morning's Q565 artifact closed with one open conjecture
(\texttt{conj:transversal}): for $\bftheta=(0,1,\dots,r-1)$ and $M$ reciprocal-free
consisting of $\lambda$ negation pairs $\{v_i,-v_i\}$ together with $f$ unpaired
elements, $\PsiC_{(1^r),\bftheta}(M)=0$ \emph{if and only if} $r\ge f+2$; it was
measured $26/26$ cells in both directions and the sign bookkeeping was the missing
step. Both halves turn out to be wrong in instructive ways, and the repaired statement
is stronger than the conjecture.

The ``if'' direction is now \textbf{proved} (Theorem~\ref{thm:main}), for a scope much
wider than $\bftheta=(0,\dots,r-1)$ and with a threshold one step lower than
conjectured on half of that scope: whenever the $\rho$-shifts $m_j=\theta_j+j-1$ are
\emph{all of the same parity} $\delta$, the vanishing holds for $r\ge f+2-\delta$. The proof is not a sign-reversing
involution: both natural candidates fail, the single-pair flip because it has fixed
points as soon as $r\le\lambda+f-1$, and the global flip because its sign is
$+1$ whenever $\binom{r}{2}$ is even. Instead the $2^{|K|}$ terms of each orbit of the pair-flipping
group $(\ZZ/2)^{K}$ are shown to sum to a rational function that is \emph{simultaneously}
divisible by $\prod_{\alpha<\beta}(x_\alpha^2-x_\beta^2)$ and of too low a degree to be
(Proposition~\ref{prop:core}). The threshold is exactly where that squeeze closes; and
when all the $m_j$ are odd the orbit sum is in addition odd in each variable, which
supplies one more factor and lowers the threshold by one. The $\delta=1$ half was
invisible to the original sweep, which only ever used $\bftheta=(0,\dots,r-1)$, i.e.\
$m_j=2(j-1)$: all even.

The ``only if'' direction is \textbf{false}, and the reason is a category error rather
than a missing case: $r\ge f+2$ is the threshold for vanishing \emph{identically on the
stratum}, and the conjecture read it as a pointwise criterion. Four exact rational
witnesses (Theorem~\ref{thm:refute}) have $r=1<f+2$ and $\PsiC=0$ with every term
nonzero; by Theorem~C of the morning artifact their locus is $\prod_t w_t=\pm1$, a
codimension-one subvariety \emph{of the stratum}, which a sampler drawing random points
of the stratum can never meet. One of the four lies in a cell the $26/26$ sweep
actually visited.
\end{abstract}

\section{What was conjectured, and the shape of the error}

\begin{conjecture}[\texttt{conj:transversal}, 2026-10-24 cycle 1, now settled]
\label{conj:old}
Let $\bftheta=(0,1,\dots,r-1)$ and let $M$ be admissible and reciprocal-free,
consisting of $\lambda$ disjoint negation pairs $\{v_i,-v_i\}$ together with $f$
unpaired elements. Then $\PsiC_{(1^r),\bftheta}(M)=0$ if and only if $r\ge f+2$.
\end{conjecture}

Two separate things were wrong, and they are worth naming before the mathematics
because they are the same two faults the morning session diagnosed in its own
predecessor, one level up.

\paragraph{(i) The sweep could not see the counterexamples, because it drew points.}
The $26/26$ was measured by drawing random $(v,u)$ from the stratum and evaluating.
A statement of the form ``$\PsiC$ vanishes at this point'' has two very different
causes: $\PsiC$ may vanish identically on the stratum, or the point may lie on a proper
subvariety of the stratum. Random draws see the first and miss the second with
probability $1$. So the sweep was measuring the identical-vanishing threshold --- which
it got right --- and reporting it as a pointwise criterion, which it is not. The
morning artifact's own \S1 says \emph{a sampler's support is a hypothesis whether or not
it is written as one}; the present fault is the next one along: \emph{a sampler's
\textbf{measure} is a hypothesis too, and it makes every positive-codimension locus
invisible.} A sweep over a stratum can only ever confirm identities.

\paragraph{(ii) ``The missing step is the sign'' named the wrong missing step.}
The conjecture was recorded with the remark that the involution reversing one negation
pair pairs the surviving terms and only the sign bookkeeping was missing. That
diagnosis was wrong twice over. The single-pair flip $\sigma_i$ is \emph{not}
fixed-point free: the orbit of a transversal $T$ that does not meet the pair
$\{v_i,-v_i\}$ is a fixed point, and such $T$ exist as soon as $r\le\lambda+f-1$. And
the global flip $\varepsilon\mapsto-\varepsilon$, which \emph{is} fixed-point free,
multiplies an orbit term by $(-1)^{\binom{L}{2}}$ when $f=0$ --- so it proves nothing
when $\binom{L}{2}$ is even, i.e.\ for $L\equiv0,1\pmod4$, and the cell
$\lambda=f+4=4,\ f=0,\ r=4$ is measured zero with $16$ nonzero terms. There is no
sign-reversing involution to find. The mechanism is a degree count.

\section{Setting}

\begin{convention}\label{conv:setup}
$M=\{w_1,\dots,w_n\}$ is a set of $n$ distinct nonzero elements of a field $k$ of
characteristic $0$ (\emph{admissible}). $M$ is \emph{reciprocal-free} if
$w_aw_b\neq1$ for all $a\le b$ --- in particular $w_a\neq\pm1$. For
$\bftheta\in\ZZ^r$ put $m_j=\theta_j+j-1$, assumed pairwise distinct. I use throughout
the alternant expansion of 2026-10-22 and the pair form of $B(T)$ of this morning:
\begin{align}
  \PsiC_{(1^r),\bftheta}(M) &= \sum_{T\subseteq M,\ |T|=r} B(T)\,A(T),
  &A(T)&=\frac{\det\bigl(x_i^{m_j}\bigr)_{i,j=1}^{r}}{\prod_{i<k}(x_k-x_i)},
  \label{eq:alt}\\
  B(T)&=\frac{N(T)}{D(T)},
  &N(T)&=\prod_{x\in T}(x^2-1)\!\!\prod_{\{x,y\}\subseteq T}\!\!(xy-1)
          \prod_{x\in T}\prod_{t\in M\setminus T}\!\!(xt-1),
  \label{eq:pairform}
\end{align}
with $T=\{x_1,\dots,x_r\}$ and
$D(T)=e^{r}\prod_{x\in T}x\prod_{x\in T}\prod_{t\in M\setminus T}(x-t)$, $e=\prod_{t\in M}t$.
\end{convention}

\begin{convention}[the negation stratum]\label{conv:stratum}
Given $\lambda,f\ge0$ let $\Sigma_{\lambda,f}\subseteq\mathbb{A}^{\lambda+f}$ be the set
of $(v_1,\dots,v_\lambda,u_1,\dots,u_f)$ for which
\[
  M(v,u)\;=\;\{v_1,-v_1,\dots,v_\lambda,-v_\lambda\}\;\cup\;\{u_1,\dots,u_f\}
\]
is admissible and has \emph{exactly} $\lambda$ negation pairs, i.e.\ $u_p\neq-u_q$ for
all $p,q$ and $u_p\neq\pm v_i$ for all $p,i$. Thus $n=2\lambda+f$. Write
$\Sigma^{\circ}_{\lambda,f}$ for the reciprocal-free part. Both are nonempty Zariski-open
subsets of $\mathbb{A}^{\lambda+f}$, hence irreducible and dense in it.
\end{convention}

\section{The parity hypothesis, and what it buys}

\begin{lemma}[alternant parity split]\label{lem:parity}
Suppose $m_j\equiv\delta\pmod2$ for all $j$, $\delta\in\{0,1\}$, and write
$m_j=2k_j+\delta$. Then for $T=\{x_1,\dots,x_r\}$,
\[
  A(T)\;=\;\Bigl(\prod_{i=1}^{r}x_i\Bigr)^{\!\delta}\cdot
           s_\mu\bigl(x_1^2,\dots,x_r^2\bigr)\cdot\!\!\prod_{1\le i<i'\le r}\!\!(x_i+x_{i'}),
\]
where $s_\mu(y)=\det(y_i^{k_j})/\prod_{i<i'}(y_{i'}-y_i)$ is the Schur function of the
partition $\mu$ with $\mu_{r+1-j}=k_j-(j-1)$.
\end{lemma}

\begin{proof}
$\det(x_i^{2k_j+\delta})=\bigl(\prod_i x_i^{\delta}\bigr)\det\bigl((x_i^2)^{k_j}\bigr)$,
by pulling $x_i^{\delta}$ out of row $i$. Also
$\prod_{i<i'}(x_{i'}^2-x_i^2)=\prod_{i<i'}(x_{i'}-x_i)\prod_{i<i'}(x_{i'}+x_i)$. Divide
the first display by $\prod_{i<i'}(x_{i'}-x_i)$ and insert the second.
\end{proof}

Verified $216/216$ exactly, $r\le4$, random rational $x$ with distinct
$|x_i|$ (\texttt{t7\_parity.py}).

\begin{corollary}\label{cor:dead}
Under the same parity hypothesis, $A(T)=0$ whenever $T$ contains both members of a
negation pair.
\end{corollary}

\begin{proof}
If $x,-x\in T$ then the corresponding rows of $(x_i^{m_j})$ are
$(x^{m_j})_j$ and $((-1)^{m_j}x^{m_j})_j=(-1)^{\delta}(x^{m_j})_j$: proportional, so
the determinant vanishes. (No appeal to Lemma~\ref{lem:parity} is needed.)
\end{proof}

\begin{convention}[orbits]\label{conv:orbits}
Fix $(v,u)\in\Sigma_{\lambda,f}$ and $r\ge1$. By Corollary~\ref{cor:dead} only the
\emph{transversal} $r$-subsets contribute to \eqref{eq:alt}, and these are indexed by
triples $(K,J,\varepsilon)$ with $K\subseteq[\lambda]$, $J\subseteq[f]$,
$|K|+|J|=r$ and $\varepsilon\in\{\pm1\}^{K}$:
\[
  T(K,J,\varepsilon)=\{a_i:=\varepsilon_iv_i\;:\;i\in K\}\cup\{u_j:j\in J\}.
\]
For fixed $(K,J)$ the $2^{|K|}$ subsets $T(K,J,\cdot)$ form one orbit of the
pair-flipping group $(\ZZ/2)^{K}\subseteq\mathrm{Sym}(M)$. Put
\[
  L:=|K|,\qquad s:=f-|J|,\qquad N:=L+s .
\]
Since $|K|+|J|=r$ we have, for \emph{every} orbit,
\begin{equation}\label{eq:Lms}
  L-s \;=\; |K|-f+|J| \;=\; r-f .
\end{equation}
\end{convention}

Equation~\eqref{eq:Lms} is the whole reason the threshold is $r\ge f+2$: it says
$r\ge f+2$ holds if and only if $L\ge s+2$ for every orbit simultaneously.

\section{The orbit factorisation}

Write $R(a,b):=(ab-1)(a+b)$, and for $\varepsilon\in\{\pm1\}^K$ write
$\eta(\varepsilon):=\prod_{i\in K}\varepsilon_i$.

\begin{lemma}[orbit factorisation]\label{lem:bridge}
Let $(v,u)\in\Sigma^{\circ}_{\lambda,f}$, let all $m_j$ be of parity $\delta$, and fix
$(K,J)$ as in Convention~\ref{conv:orbits}. Then there is a constant
$C_{K,J}\in k^{\times}$, independent of $\varepsilon$, with
\begin{equation}\label{eq:bridge}
  B\bigl(T(K,J,\varepsilon)\bigr)\,A\bigl(T(K,J,\varepsilon)\bigr)
  \;=\; C_{K,J}\cdot\eta(\varepsilon)^{\delta}\cdot
  \prod_{\substack{i<i'\\ i,i'\in K}}\frac{1}{R(a_i,a_{i'})}\cdot
  \prod_{i\in K}\ \prod_{p\in[f]\setminus J}\frac{1}{R(a_i,-u_p)} .
\end{equation}
\end{lemma}

\begin{proof}
Write $T=T(K,J,\varepsilon)$ and $S=M\setminus T$, so that
\[
  S=\{-a_i:i\in K\}\ \sqcup\ \{v_k,-v_k:k\in[\lambda]\setminus K\}
      \ \sqcup\ \{u_p:p\in[f]\setminus J\}.
\]
Throughout, $\doteq$ means ``equal up to a nonzero factor not depending on
$\varepsilon$''. Flipping $\varepsilon_i$ replaces $a_i$ by $-a_i$ and fixes $a_i^2$;
so any expression in the $a_i^2$, the $v_k$ and the $u_p$ is $\varepsilon$-free.
I collect the factors of $N(T)A(T)/D(T)$ by the pair of elements they involve.

\smallskip\noindent\textbf{Singletons.} $\prod_{x\in T}(x^2-1)$ from $N(T)$ is
$\varepsilon$-free. $s_\mu(T^2)$ from Lemma~\ref{lem:parity} is $\varepsilon$-free,
because $T^2=\{v_i^2\}_{i\in K}\cup\{u_j^2\}_{j\in J}$. The factor
$(\prod_{x\in T}x)^{\delta}$ of Lemma~\ref{lem:parity} contributes
$\eta(\varepsilon)^{\delta}$, and the factor $\prod_{x\in T}x$ in $D(T)$ contributes
$\eta(\varepsilon)^{-1}$; the sub-product of $\prod_{x\in T}\prod_{t\in S}(x-t)$ with
$x=a_i,\ t=-a_i$ is $\prod_{i\in K}2a_i\doteq\eta(\varepsilon)$, contributing a second
$\eta(\varepsilon)^{-1}$. As $\eta=\pm1$, these three give
$\eta^{\delta-2}=\eta^{\delta}$. The sub-product with $x=a_i,\ t=-a_i$ of
$\prod_{x\in T}\prod_{t\in S}(xt-1)$ is $\prod_i(-a_i^2-1)$, $\varepsilon$-free.

\smallskip\noindent\textbf{Pairs $i<i'$ in $K$.} Put $P=a_ia_{i'}$, so $P^2$ is
$\varepsilon$-free. The contributing factors are: $(a_i+a_{i'})$ from
Lemma~\ref{lem:parity}; $(P-1)$ from $\prod_{\{x,y\}\subseteq T}(xy-1)$;
$(-P-1)^2$ from $\prod_{x\in T}\prod_{t\in S}(xt-1)$ at $(x,t)=(a_i,-a_{i'})$ and
$(a_{i'},-a_i)$; and $(a_i+a_{i'})^2$ in $D(T)$ from
$\prod_{x\in T}\prod_{t\in S}(x-t)$ at the same two ordered pairs. The net factor is
\[
  \frac{(a_i+a_{i'})\,(P-1)\,(P+1)^2}{(a_i+a_{i'})^{2}}
  \;=\;\frac{(P^2-1)(P+1)}{a_i+a_{i'}}
  \;\doteq\;\frac{P+1}{a_i+a_{i'}}
  \;=\;\frac{P^2-1}{(P-1)(a_i+a_{i'})}
  \;\doteq\;\frac{1}{R(a_i,a_{i'})}.
\]

\smallskip\noindent\textbf{Pairs $i\in K$, $j\in J$.} The factors are $(a_i+u_j)$
(Lemma~\ref{lem:parity}), $(a_iu_j-1)$ (from $\prod_{\{x,y\}\subseteq T}$),
$(-a_iu_j-1)$ (from $\prod_{x\in T}\prod_{t\in S}(xt-1)$ at $x=u_j,t=-a_i$), and
$(u_j+a_i)$ in $D(T)$ (from $\prod_{x\in T}\prod_{t\in S}(x-t)$ at $x=u_j,t=-a_i$). The
net factor is $-(a_i^2u_j^2-1)$, which is $\varepsilon$-free. (These pairs drop out
entirely; this is why $J$ enters only through its complement.)

\smallskip\noindent\textbf{Pairs $i\in K$, $p\notin J$.} The factors are
$(a_iu_p-1)$ (from $\prod_{x\in T}\prod_{t\in S}(xt-1)$) and $(a_i-u_p)$ in $D(T)$.
The net factor is
\[
  \frac{a_iu_p-1}{a_i-u_p}=\frac{a_i^2u_p^2-1}{(a_iu_p+1)(a_i-u_p)}
  \doteq\frac{1}{(a_iu_p+1)(a_i-u_p)}
  =\frac{-1}{\bigl(a_i(-u_p)-1\bigr)\bigl(a_i+(-u_p)\bigr)}
  \doteq\frac{1}{R(a_i,-u_p)} .
\]

\smallskip\noindent\textbf{Pairs $i\in K$, $k\notin K$.} Both $v_k$ and $-v_k$ lie in
$S$, so these factors occur in $\pm$ pairs:
$(a_iv_k-1)(-a_iv_k-1)=-(a_i^2v_k^2-1)$ and $(a_i-v_k)(a_i+v_k)=a_i^2-v_k^2$, both
$\varepsilon$-free. \emph{Untouched pairs contribute nothing}; this is why $\lambda$ does
not appear in the final criterion.

\smallskip\noindent\textbf{Everything else} ($e^r$; pairs inside $J$; pairs
$j\in J$ with $k\notin K$ or $p\notin J$) involves no $a_i$ and is $\varepsilon$-free.

Multiplying the four net contributions gives \eqref{eq:bridge}. Each factor of
$C_{K,J}$ is a difference or product of distinct elements of $M$, or of the form
$w_aw_b-1$, and so is nonzero because $M$ is admissible and reciprocal-free; hence
$C_{K,J}\neq0$.
\end{proof}

Verified exactly, $63/63$ orbits over $10$ cells with $\lambda\le4$, $f\le3$
(\texttt{t4\_bridge.py}): the quotient of $B(T)A(T)$ by the displayed product takes a
single value on each orbit.

\begin{lemma}[linearisation]\label{lem:lin}
Put $\sigma(x)=x-x^{-1}$. Then $R(a,b)=ab\,\bigl(\sigma(a)+\sigma(b)\bigr)$, and
$\sigma(-x)=-\sigma(x)$.
\end{lemma}

\begin{proof}
$ab(\sigma(a)+\sigma(b))=ab\bigl(a-a^{-1}+b-b^{-1}\bigr)=a^2b+ab^2-b-a
 =(ab-1)(a+b)$.
\end{proof}

\begin{corollary}[the orbit sum as a sign-twisted Cauchy sum]\label{cor:orbitsum}
In the situation of Lemma~\ref{lem:bridge}, set
\[
  x_i=\sigma(v_i)\ (i\in K),\qquad c_p=\sigma(-u_p)=-\sigma(u_p)\ (p\in[f]\setminus J),
  \qquad \nu=N-1+\delta .
\]
Then, for a nonzero constant $C'_{K,J}$ independent of $\varepsilon$,
\begin{equation}\label{eq:orbitsum}
  \sum_{\varepsilon\in\{\pm1\}^{K}}\!\! B(T(K,J,\varepsilon))A(T(K,J,\varepsilon))
  \;=\;C'_{K,J}\!\!\sum_{\varepsilon\in\{\pm1\}^{L}}\!\!
     \eta(\varepsilon)^{\nu}
     \prod_{\alpha<\beta\le L}\!\frac{1}{\varepsilon_\alpha x_\alpha+\varepsilon_\beta x_\beta}
     \prod_{\alpha\le L}\prod_{p\le s}\!\frac{1}{\varepsilon_\alpha x_\alpha+c_p}.
\end{equation}
\end{corollary}

\begin{proof}
Form the graph $G$ on the $N$ vertices $K\sqcup([f]\setminus J)$, with an edge for every
pair having at least one endpoint in $K$; assign the value $a_i$ to $i\in K$ and $-u_p$
to $p\notin J$. The product in \eqref{eq:bridge} is $\prod_{\text{edges}}R^{-1}$. By
Lemma~\ref{lem:lin},
$\prod_{\text{edges}}R=\prod_{\text{vertices}}(\text{value})^{\deg}\cdot
 \prod_{\text{edges}}(\sigma_\alpha+\sigma_\beta)$,
and $\deg=N-1$ at a vertex of $K$ and $\deg=L$ at a vertex of $[f]\setminus J$. Since
$\sigma(\varepsilon_iv_i)=\varepsilon_i\sigma(v_i)=\varepsilon_ix_i$ and
$\prod_{i\in K}a_i^{N-1}=\eta^{N-1}\prod_i v_i^{N-1}$, the right-hand side of
\eqref{eq:bridge} equals a constant times
$\eta^{\delta}\eta^{-(N-1)}\prod_{\text{edges}}(\varepsilon_\alpha x_\alpha
+\varepsilon_\beta x_\beta)^{-1}$, and $\eta^{\delta-(N-1)}=\eta^{\nu}$ as $\eta=\pm1$.
Summing over $\varepsilon$ gives \eqref{eq:orbitsum}. The denominators are nonzero:
$\varepsilon_\alpha x_\alpha+\varepsilon_\beta x_\beta=R(a_\alpha,a_\beta)/(a_\alpha
a_\beta)$ is nonzero because $a_\alpha a_{\beta}\neq1$ (reciprocal-freeness) and
$a_\alpha+a_\beta\neq0$ ($\alpha,\beta$ lie in different negation pairs); and
$\varepsilon_\alpha x_\alpha+c_p=R(a_\alpha,-u_p)/(-a_\alpha u_p)$ is nonzero because
$a_\alpha u_p\neq-1$ (apply reciprocal-freeness to $-a_\alpha\in M$ and $u_p$) and
$a_\alpha\neq u_p$ (distinctness).
\end{proof}

\section{The core identity}

\begin{proposition}[core identity]\label{prop:core}
Let $L,s\ge0$, $N=L+s$, and $\nu\in\ZZ$. Let $x_1,\dots,x_L,c_1,\dots,c_s$ be
independent indeterminates over $\QQ$ and set
\[
  Z \;=\; \sum_{\varepsilon\in\{\pm1\}^{L}}
     \Bigl(\prod_{\alpha=1}^{L}\varepsilon_\alpha\Bigr)^{\!\nu}
     \prod_{1\le\alpha<\beta\le L}\frac{1}{\varepsilon_\alpha x_\alpha+\varepsilon_\beta x_\beta}
     \prod_{\alpha=1}^{L}\prod_{p=1}^{s}\frac{1}{\varepsilon_\alpha x_\alpha+c_p}.
\]
Then $Z=0$ in each of the two cases
\[
  \text{(i) } L\ge s+2;\qquad\qquad \text{(ii) } L\ge s+1\ \text{and}\ \nu\ \text{odd}.
\]
\end{proposition}

\begin{lemma}[parity of $Z$ in each free variable]\label{lem:Zodd}
With the notation of Proposition~\ref{prop:core},
$Z|_{x_\alpha\mapsto-x_\alpha}=(-1)^{\nu}Z$ for every $\alpha\le L$.
\end{lemma}

\begin{proof}
Take $\alpha=1$. Substituting $-x_1$ for $x_1$ in the summand indexed by $\varepsilon$
replaces every occurrence of $\varepsilon_1x_1$ by $(-\varepsilon_1)x_1$ and touches
nothing else, so that summand becomes
$\eta(\varepsilon)^{\nu}$ times the $\varepsilon$-independent-of-$x_1$ part of the
summand indexed by $\varepsilon'$, where $\varepsilon'_1=-\varepsilon_1$ and
$\varepsilon'_\beta=\varepsilon_\beta$ for $\beta>1$. Since
$\eta(\varepsilon)=-\eta(\varepsilon')$, that summand equals $(-1)^{\nu}$ times the
$\varepsilon'$-summand. As $\varepsilon\mapsto\varepsilon'$ is a bijection of
$\{\pm1\}^L$, summing gives $Z(-x_1)=(-1)^{\nu}Z(x_1)$.
\end{proof}

\begin{proof}
Put
\[
  D\;=\;\prod_{1\le\alpha<\beta\le L}\bigl(x_\alpha^2-x_\beta^2\bigr)\cdot
        \prod_{\alpha=1}^{L}\prod_{p=1}^{s}\bigl(c_p^2-x_\alpha^2\bigr),
  \qquad
  \Delta\;=\;\prod_{1\le\alpha<\beta\le L}\bigl(x_\alpha^2-x_\beta^2\bigr).
\]
Each summand's denominator is a product of distinct irreducible linear forms, each
dividing $D$ exactly once: $\varepsilon_\alpha x_\alpha+\varepsilon_\beta x_\beta$ is
$\pm(x_\alpha\pm x_\beta)$, and $\varepsilon_\alpha x_\alpha+c_p$ is
$c_p\pm x_\alpha$. Hence $P:=ZD$ is a polynomial, and it suffices to show $P=0$.

\smallskip\noindent\textbf{Step 1 (degree in $x_1$).} Exactly $N-1$ factors of each
summand involve $x_1$, namely the $L-1$ factors
$\varepsilon_1x_1+\varepsilon_\beta x_\beta$ $(2\le\beta\le L)$ and the $s$ factors
$\varepsilon_1x_1+c_p$; each is of degree $1$ in $x_1$, and no numerator involves
$x_1$. So each summand is $O\bigl(x_1^{-(N-1)}\bigr)$ as $x_1\to\infty$, and therefore
so is $Z$. Since $\deg_{x_1}D=2(L-1)+2s$,
\begin{equation}\label{eq:degbound}
  \deg_{x_1}P\;\le\;2(L-1)+2s-(N-1)\;=\;L+s-1\;=\;N-1 .
\end{equation}

\smallskip\noindent\textbf{Step 2 (divisibility by $\Delta$).} Fix $\alpha<\beta\le L$
and $\varsigma\in\{\pm1\}$, and let $\ell=x_\alpha-\varsigma x_\beta$, an irreducible
linear form, with $V=\{\ell=0\}$. For a summand indexed by $\varepsilon$, the factor
$\varepsilon_\alpha x_\alpha+\varepsilon_\beta x_\beta$ vanishes on $V$ precisely when
$\varepsilon_\alpha\varsigma+\varepsilon_\beta=0$, i.e.\ when
$\varepsilon_\alpha\varepsilon_\beta=-\varsigma$; no other factor vanishes at a generic
point of $V$. Let $\iota$ flip both $\varepsilon_\alpha$ and $\varepsilon_\beta$. Then
$\iota$ preserves the set of singular indices (it preserves
$\varepsilon_\alpha\varepsilon_\beta$) and is fixed-point free on it, so the singular
summands are partitioned into $\iota$-pairs. Write a singular summand as
$T(\varepsilon)=g_\varepsilon/(\varepsilon_\alpha x_\alpha+\varepsilon_\beta x_\beta)$
with $g_\varepsilon$ regular at the generic point of $V$. Put
$A=\varepsilon_\alpha x_\alpha$ and $B=\varepsilon_\beta x_\beta$, so $B=-A$ on $V$.
Comparing $g_\varepsilon$ with $g_{\iota\varepsilon}$ factor by factor:
\begin{itemize}
\item $\prod\varepsilon$ is unchanged by $\iota$ (two flips), so the sign
      $\eta^{\nu}$ is unchanged;
\item for $\gamma\notin\{\alpha,\beta\}$, $\gamma\le L$, the pair of factors
      $(A+X)(B+X)$ with $X=\varepsilon_\gamma x_\gamma$ becomes $(-A+X)(-B+X)$; on $V$
      both equal $X^2-A^2$;
\item likewise $(A+c_p)(B+c_p)$ becomes $(-A+c_p)(-B+c_p)$, both equal to $c_p^2-A^2$
      on $V$;
\item factors involving neither $\alpha$ nor $\beta$ are unchanged.
\end{itemize}
Hence $g_\varepsilon=g_{\iota\varepsilon}$ on $V$, and since
$\varepsilon_\alpha x_\alpha+\varepsilon_\beta x_\beta=\pm\ell$,
\[
  T(\varepsilon)+T(\iota\varepsilon)
  =\frac{g_\varepsilon-g_{\iota\varepsilon}}{\pm\ell}
\]
is regular at the generic point of $V$, because $\ell$ divides the numerator. Summing
over the $\iota$-pairs, $Z$ is regular on $V$. As $D$ vanishes to order exactly one on
$V$, $P=ZD$ vanishes on $V$, i.e.\ $\ell\mid P$. Running over all $\alpha<\beta\le L$
and both $\varsigma$ gives $2\binom{L}{2}$ pairwise non-associate irreducible linear
forms dividing $P$, whence
\[
  \Delta\;=\;\prod_{\alpha<\beta\le L}(x_\alpha-x_\beta)(x_\alpha+x_\beta)\ \Bigm|\ P .
\]

\smallskip\noindent\textbf{Step 3 (divisibility by $\prod_\alpha x_\alpha$ when $\nu$
is odd).} $D$ is even in each $x_\alpha$, so by Lemma~\ref{lem:Zodd} $P=ZD$ satisfies
$P|_{x_\alpha\mapsto-x_\alpha}=(-1)^{\nu}P$. If $\nu$ is odd then $P$ is odd in each
$x_\alpha$, so $P|_{x_\alpha=0}=0$ and $x_\alpha\mid P$. No $x_\gamma$ divides any
$x_\alpha\pm x_\beta$, so $x_1,\dots,x_L$ together with the $2\binom L2$ forms of Step 2
are pairwise non-associate irreducibles, and
\[
  \Delta\cdot\textstyle\prod_{\alpha=1}^{L}x_\alpha\ \Bigm|\ P
  \qquad(\nu\ \text{odd}).
\]

\smallskip\noindent\textbf{Step 4 (the squeeze).} $\deg_{x_1}\Delta=2(L-1)$, since
exactly the $L-1$ factors $x_1^2-x_\beta^2$ involve $x_1$. Suppose $P\neq0$.

In case (i), $\Delta\mid P$ forces $\deg_{x_1}P\ge2(L-1)$, so with \eqref{eq:degbound}
\[
  2(L-1)\;\le\;N-1\;=\;L+s-1,\qquad\text{i.e.}\qquad L\le s+1 ,
\]
contradicting $L\ge s+2$.

In case (ii), $\nu$ is odd, so Step 3 gives $\deg_{x_1}P\ge2(L-1)+1$ and
\eqref{eq:degbound} gives
\[
  2(L-1)+1\;\le\;L+s-1,\qquad\text{i.e.}\qquad L\le s ,
\]
contradicting $L\ge s+1$. In both cases $P=0$, hence $Z=0$.
\end{proof}

\begin{remark}\label{rem:sharp}
Both squeezes are exactly sharp: $2(L-1)>L+s-1\iff L\ge s+2$ and
$2(L-1)+1>L+s-1\iff L\ge s+1$, so no slack is discarded and no case is left over. Note
that Step 2 is \emph{not} a sign-reversing involution on the sum --- $\iota$ reverses
signs only on the hyperplane $V$, and only on the singular summands. The involution
organises a residue cancellation, not the sum.

I found case (ii) by running \texttt{t10\_proofinputs.py} as a control on
case (i) and reading its four disagreements rather than its $40$ agreements: every one
was $L=s+1$ with $\nu$ odd. \emph{The control that fires is the one carrying the
information.}
\end{remark}

\begin{remark}
Proposition~\ref{prop:core} says nothing outside (i) and (ii), and indeed $Z\neq0$
there: of the $44$ cells with $L\le5$, $s\le4$, $N\le7$ and $\nu=N-1+\delta$,
$\delta\in\{0,1\}$, exactly the $20$ satisfying (i) or (ii) have $Z=0$, measured at $12$
random rational points each (\texttt{t10\_proofinputs.py}).
\end{remark}

\section{The theorem}

\begin{theorem}[transversal cancellation]\label{thm:main}
Let $M$ be admissible with exactly $\lambda$ negation pairs and $f=n-2\lambda$ unpaired
elements, and let $\bftheta\in\ZZ^{r}$ have $\rho$-shifts $m_j=\theta_j+j-1$ pairwise
distinct and \emph{all of the same parity} $\delta\in\{0,1\}$. If
\[
  r\;\ge\;f+2-\delta
\]
then $\PsiC_{(1^r),\bftheta}(M)=0$. (So $r\ge f+2$ when all $m_j$ are even, and the one
step better $r\ge f+1$ when all $m_j$ are odd.)
\end{theorem}

\begin{proof}
First assume $M$ reciprocal-free, i.e.\ $(v,u)\in\Sigma^{\circ}_{\lambda,f}$. By
Corollary~\ref{cor:dead} the non-transversal $r$-subsets contribute $0$ to
\eqref{eq:alt}, so
\[
  \PsiC_{(1^r),\bftheta}(M)
  =\sum_{(K,J)}\ \sum_{\varepsilon\in\{\pm1\}^{K}}
     B\bigl(T(K,J,\varepsilon)\bigr)A\bigl(T(K,J,\varepsilon)\bigr),
\]
the outer sum over $K\subseteq[\lambda]$, $J\subseteq[f]$ with $|K|+|J|=r$. Fix one
$(K,J)$ and let $L=|K|$, $s=f-|J|$, $N=L+s$. By Corollary~\ref{cor:orbitsum} the inner
sum is $C'_{K,J}$ times the specialisation of the rational function $Z$ of
Proposition~\ref{prop:core} at $x_\alpha=\sigma(v_\alpha)$, $c_p=-\sigma(u_p)$,
$\nu=N-1+\delta$; that specialisation is legitimate because no denominator vanishes
there (end of the proof of Corollary~\ref{cor:orbitsum}). It remains to check that
Proposition~\ref{prop:core} applies. By \eqref{eq:Lms}, $L-s=r-f$ for this and every
orbit.
\begin{itemize}
\item If $r\ge f+2$ then $L\ge s+2$: case (i).
\item If $r=f+1$ (which is the remaining possibility only when $\delta=1$) then
      $L=s+1$, so $N=2s+1$ and $\nu=N-1+\delta=2s+\delta$ is odd: case (ii).
\end{itemize}
Either way $Z=0$ identically, so the inner sum is $0$. Every orbit contributes $0$, so
$\PsiC=0$.

Now let $M$ be merely admissible. By \eqref{eq:alt} and \eqref{eq:pairform},
$\PsiC_{(1^r),\bftheta}$ is a rational function of $(v,u)$ whose denominators are
products of $e^r$, of differences $x-t$ of distinct elements of $M$, and of
$\prod_{i<k}(x_k-x_i)$; all are nonzero on $\Sigma_{\lambda,f}$, so $\PsiC$ is regular
there. $\Sigma_{\lambda,f}$ is a nonempty open subset of
$\mathbb{A}^{\lambda+f}$, hence irreducible, and $\Sigma^{\circ}_{\lambda,f}$ is a
nonempty open subset of it, hence dense. A regular function vanishing on a dense subset
of an irreducible variety vanishes identically. So $\PsiC=0$ on all of
$\Sigma_{\lambda,f}$.
\end{proof}

\begin{corollary}\label{cor:scope}
$\PsiC_{(1^r),\bftheta}(M)=0$ for every $\bftheta=(c,c+1,\dots,c+r-1)$ with $c\in\ZZ$
(then $m_j=2(j-1)+c$, of parity $\delta\equiv c$) as soon as $r\ge f+2-(c\bmod2)$; in
particular for $\bftheta=(0,1,\dots,r-1)$ and $r\ge f+2$, which is the ``if'' half of
Conjecture~\ref{conj:old}.
\end{corollary}

\begin{remark}[the sharpened threshold is new, and the old sweep could not have seen it]
The $26/26$ sweep behind Conjecture~\ref{conj:old} used only
$\bftheta=(0,1,\dots,r-1)$, i.e.\ $m_j=2(j-1)$, which is $\delta=0$ throughout. It
therefore reported $r=f+1$ as the clean boundary (``$r=f+1$ gave $0$ zeros in every
cell''), and that is correct --- \emph{for even $m$}. For odd $m$ the boundary is one
step lower. Measured: all $m_j$ odd at $r=f+1$ gives $\PsiC=0$ in $160/160$ draws
($\lambda\le5$, $f\le3$, both engines agreeing on every draw), while the two controls
stay green --- all $m_j$ even at $r=f+1$ is nonzero $160/160$, and all $m_j$ odd at
$r=f$ (one below the new threshold) is nonzero $120/120$ (\texttt{t11\_sharp.py}). The
full criterion $r\ge f+2-\delta$ gives $0$ in $368/368$ swept draws.
\end{remark}

Measured: $170/170$ draws zero for same-parity $m$ of arbitrary spacing, $\lambda\le4$,
$f\le2$, $r\ge f+2$, $\delta\in\{0,1\}$; and $126/126$ for the shifted arithmetic
progressions $\bftheta=(c,\dots,c+r-1)$, $c\in[-3,3]$ (\texttt{t6\_scope.py},
\texttt{t7\_parity.py}). The grid of \texttt{t2\_grid.py} agrees in $89/89$ cells,
$\lambda\le4$, $f\le3$, $r\le6$.

\begin{remark}[the parity hypothesis is load-bearing, and not necessary]
\label{rem:parity}
A control run on $m$ of \emph{mixed} parity gives $149/170$ nonzero
(\texttt{t7\_parity.py}), so the hypothesis is doing real work. It is not necessary,
however: see \S\ref{sec:open}.
\end{remark}

\section{The converse is false}

\begin{theorem}\label{thm:refute}
The ``only if'' half of Conjecture~\ref{conj:old} is false. Each of the following $M$ is
admissible and reciprocal-free, has $r<f+2$, has every term $B(T)A(T)\neq0$, and has
$\PsiC_{(1^r),(0,\dots,r-1)}(M)=0$:
\[
\begin{array}{llll}
  \text{(a)} & M=\{2,-2,3,-3,\tfrac16,-\tfrac16\}, & \lambda=3,\ f=0,\ r=1, & \textstyle\prod_t w_t=-1;\\
  \text{(b)} & M=\{2,-\tfrac12\}, & \lambda=0,\ f=2,\ r=1, & \textstyle\prod_t w_t=-1;\\
  \text{(c)} & M=\{3,-3,\tfrac19\}, & \lambda=1,\ f=1,\ r=1, & \textstyle\prod_t w_t=-1;\\
  \text{(d)} & M=\{2,-2,5,-5,\tfrac1{100}\}, & \lambda=2,\ f=1,\ r=1, & \textstyle\prod_t w_t=+1.
\end{array}
\]
\end{theorem}

\begin{proof}
Each is a finite exact computation, carried out in \texttt{fractions.Fraction} by two
independent implementations of $\PsiC$ --- the head-injection sum and the alternant
expansion \eqref{eq:alt} --- which agree on every one (\texttt{t1\_refute.py}).
Conceptually, all four are instances of Theorem~C of the morning artifact
(\texttt{thm:r1}) at $r=1$, $\bftheta=(0)$, which gives
\begin{equation}\label{eq:r1}
  \PsiC_{(1),(0)}(M)\;=\;1-\Bigl(\prod_{t\in M}t\Bigr)^{-2}
\end{equation}
for every admissible $M$; so the zero locus at $r=1$, $\bftheta=(0)$ is exactly
$\{\prod_t w_t=\pm1\}$, and each of (a)--(d) is a reciprocal-free point of the relevant
stratum lying on it.
\end{proof}

\begin{remark}[where the stratum meets $\prod_t w_t=\pm1$]\label{rem:meet}
On $\Sigma^{\circ}_{\lambda,f}$ one has $\prod_t w_t=(-1)^{\lambda}
\bigl(\prod_i v_i\bigr)^{2}\prod_p u_p$, so $\{\prod_t w_t=\pm1\}$ meets
$\Sigma^{\circ}_{\lambda,f}$ in a nonempty hypersurface whenever $f\ge1$ (solve for
$u_f$), and for $f=0$ whenever $\lambda\ge3$ (example (a)). It does \emph{not} meet it
for $f=0$ and $\lambda\le2$: $\lambda=1$ needs $v_1^2=\pm1$ and $\lambda=2$ needs
$v_1v_2=\pm1$, and all of $v_1=\pm1$, $v_1v_2=1$, $v_1(-v_2)=1$ are excluded by
reciprocal-freeness. So the smallest counterexamples are exactly where they are.
\end{remark}

\begin{remark}[one of the four was inside the swept domain]
The sweep behind Conjecture~\ref{conj:old} declared the domain $\lambda\le4$, $f\le3$,
$n\le7$, $r\in\{f+1,f+2,f+3\}$, $10$ draws per cell. Witness (a) has
$\lambda=3$, $f=0$, $n=6$, $r=1=f+1$: that cell \emph{was} visited, and reported
``$0$ zeros''. Nothing was wrong with the draws. $\{\prod_tw_t=\pm1\}$ has
codimension~$1$ in the stratum, and a random rational point misses it almost surely.
\end{remark}

\begin{corollary}[the corrected statement]\label{cor:corrected}
Let $\bftheta$ have pairwise distinct $\rho$-shifts, all of parity $\delta$. On
$\Sigma_{\lambda,f}$:
\begin{itemize}
\item if $r\ge f+2-\delta$ then $\PsiC_{(1^r),\bftheta}$ vanishes \emph{identically}
      (Theorem~\ref{thm:main});
\item if $r\le f+1-\delta$ then $\PsiC_{(1^r),\bftheta}$ does not vanish identically,
      but its zero locus in $\Sigma_{\lambda,f}$ is still nonempty whenever
      $\Sigma^{\circ}_{\lambda,f}$ meets it --- for $r=1$, $\bftheta=(0)$ it is exactly
      $\{\prod_tw_t=\pm1\}$ by \eqref{eq:r1}, and Remark~\ref{rem:meet} says when that
      is nonempty.
\end{itemize}
So the threshold of Conjecture~\ref{conj:old} is right as a threshold for
\emph{identical} vanishing --- after the correction $f+2\rightsquigarrow f+2-\delta$ ---
and wrong as a pointwise criterion. The second bullet's non-vanishing is \emph{computed},
not proved: $0$ zeros in $6$ independent draws in each of the $41$ cells with $r\le f+1$
of \texttt{t2\_grid.py} ($\delta=0$), plus the two green controls of
\texttt{t11\_sharp.py}.
\end{corollary}

\section{A second, \texorpdfstring{$\bftheta$}{theta}-independent mechanism}

Theorem~\ref{thm:main} needs the parity hypothesis. For $f=0$ there is a cruder
mechanism that needs nothing of $\bftheta$ at all, and it is worth recording because my
first version of it was wrong and a control caught it.

\begin{proposition}\label{prop:global}
Let $f=0$, so $M=-M$. Then for every $\bftheta$,
\[
  \PsiC_{(1^r),\bftheta}(M)\;=\;(-1)^{\,\sum_j m_j+\binom r2}\,\PsiC_{(1^r),\bftheta}(M);
  \qquad\text{hence }\PsiC=0\text{ if }\textstyle\sum_j m_j+\binom r2\text{ is odd.}
\]
\end{proposition}

\begin{proof}
$T\mapsto-T$ is a bijection of the $r$-subsets of $M$, and
$M\setminus(-T)=-(M\setminus T)$.

\emph{$N$ is invariant.} Every factor of $N(T)$ in \eqref{eq:pairform} is of the form
$x^2-1$, $xy-1$ or $xt-1$ with both arguments ranging over sets that are negated
together, so each factor is unchanged: $(-x)(-y)-1=xy-1$.

\emph{$D$ is invariant.} $e=\prod_{t\in M}t$ is unchanged ($M=-M$). The factor
$\prod_{x\in T}x$ acquires $(-1)^{r}$, and
$\prod_{x\in T}\prod_{t\in M\setminus T}(x-t)$ acquires $(-1)^{r(n-r)}$, for a total of
$(-1)^{r(n-r+1)}$. Since $f=0$ forces $n=2\lambda$ even, $r(n-r+1)\equiv r(1-r)=-r(r-1)
\equiv0\pmod 2$. So $D(-T)=D(T)$, and therefore $B(-T)=B(T)$.

\emph{$A$ picks up the sign.} Negating every $x_i$ scales column $j$ of
$(x_i^{m_j})$ by $(-1)^{m_j}$, so $\det\mapsto(-1)^{\sum_j m_j}\det$, while
$\prod_{i<k}(x_k-x_i)\mapsto(-1)^{\binom r2}\prod_{i<k}(x_k-x_i)$. Hence
$A(-T)=(-1)^{\sum_jm_j+\binom r2}A(T)$.

Reindexing \eqref{eq:alt} by $T\mapsto-T$ gives the claim.
\end{proof}

Verified $140/140$ (f=0, $\sum_jm_j+\binom r2$ odd, arbitrary $m$, $\lambda\le5$,
$r\le5$); the control, $f=0$ with $\sum_jm_j+\binom r2$ \emph{even} and $m$ of mixed
parity, is nonzero $122/126$ (\texttt{t9\_corrected.py}).

\begin{remark}[the control found the sign]
I first stated Proposition~\ref{prop:global} with the factor $(-1)^{\binom r2}$ alone,
having silently assumed $\sum_jm_j$ even (true for $\bftheta=(0,\dots,r-1)$, where
$m_j=2(j-1)$, and false in general). Run as a prediction for arbitrary $\bftheta$ it
scored $26/60$ --- about half, which is the signature of an unmodelled sign. Inserting
$(-1)^{\sum_jm_j}$ gives $140/140$. \emph{A sign specialised at the $\bftheta$ I had in
mind read as a sign that was not there.}
\end{remark}

\section{What is still open}\label{sec:open}

\begin{remark}[the parity hypothesis is not necessary]
Scanning all $16$ parity patterns of $m$ at $\lambda=3$, $f=1$, $r=4$ (so $r\ge f+2$),
$6$ independent draws per pattern, identical vanishing occurs for exactly six patterns
(\texttt{t9\_corrected.py}):
\[
  (0,0,0,0),\quad (1,1,1,1),\quad (0,1,1,1),\quad(1,0,1,1),\quad(1,1,0,1),\quad(1,1,1,0),
\]
i.e.\ for $\#\{j:m_j\text{ even}\}\in\{0,1,4\}$ and not for $\{2,3\}$. The first two are
Theorem~\ref{thm:main}. The other four are a genuine further mechanism: all $6$ draws
give $0$, so the vanishing is identical on the stratum, not accidental. The Laplace
expansion of $\det(x_i^{m_j})$ along the even and odd columns writes $A(T)$ as a sum
over $\binom{r}{\#\text{even}}$ terms, each carrying a \emph{different} character of
$(\ZZ/2)^{K}$, which is precisely why Lemma~\ref{lem:bridge} does not apply --- and
presumably why one even column behaves like none. I have not pursued it.
\end{remark}

Both $\#\text{even}=0$ and $\#\text{even}=r$ are Theorem~\ref{thm:main} (with
$\delta=1$ and $\delta=0$), and at this cell ($r=4\ge f+2=3$) the sharpened threshold
buys nothing extra, so the four patterns with $\#\text{even}=1$ are genuinely outside
everything proved here.

This is \textbf{Q578}: \emph{for which parity profiles of $(m_1,\dots,m_r)$ does the
transversal cancellation hold, and with what threshold in $r$?} Measured answer at
$\lambda=3,f=1,r=4$: $\#\text{even}\in\{0,1,4\}$. Proved: $\#\text{even}\in\{0,r\}$,
with threshold $r\ge f+2$ for $\#\text{even}=r$ and $r\ge f+1$ for $\#\text{even}=0$.

Also still open, carried from this morning: Q575 (the zero locus on the reciprocal-free
stratum for $r\ge2$) and Q577 (the type-C contour formula). Today's Corollary
\ref{cor:corrected} sharpens Q575: on the negation stratum with $r\le f+1$, the locus is
cut out by \eqref{eq:r1} at $r=1$ and is unknown for $r\ge2$.

\section{Verification}\label{sec:verif}

All arithmetic is exact: \texttt{fractions.Fraction} or SymPy over $\QQ$; no floating
point. Two independent implementations of $\PsiC$ are used throughout --- the
head-injection sum over ordered $r$-tuples (\texttt{code-1023-q559/psi.py},
eq.~(1) of the morning artifact) and this session's
\texttt{psi\_pair} (alternant expansion with $B$ in the pair form) --- and they agree
$350/350$ over $n\le5$, $r\le4$, \emph{with signs allowed and no hygiene exclusions}
(\texttt{t0\_crosscheck.py}). That last clause matters: the Q559 sampler excluded signs,
so a cross-check inherited from it would have been blind on the axis this paper lives on.

\begin{center}
\small
\begin{tabular}{llll}
\toprule
file & what & result & domain\\
\midrule
\texttt{t0\_crosscheck.py} & engine vs.\ head sum & $350/350$ & $n\le5$, $r\le4$, signed\\
\texttt{t1\_refute.py} & Theorem~\ref{thm:refute} & 4 witnesses & exact, both engines\\
\texttt{t2\_grid.py} & $r\ge f+2\Leftrightarrow$ identical $0$ & $89/89$ cells & $\lambda\le4,f\le3,r\le6$\\
\texttt{t3\_classes.py} & every orbit sum vanishes alone & $6$ cells, all orbits & incl.\ $L=4$, $L=3,s=1$\\
\texttt{t4\_bridge.py} & Lemma~\ref{lem:bridge}, Lemma~\ref{lem:lin} & $63/63$ orbits & $\lambda\le4$, $f\le3$\\
\texttt{t6\_scope.py} & $\bftheta=(c,\dots,c+r-1)$ & $126/126$ & $c\in[-3,3]$\\
\texttt{t7\_parity.py} & Lemma~\ref{lem:parity}; same-parity $m$ & $216/216$; $170/170$ & $r\le4$; $\lambda\le4$\\
\texttt{t9\_corrected.py} & Proposition~\ref{prop:global}; parity scan & $140/140$; 16 patterns & $\lambda\le5$\\
\texttt{t10\_proofinputs.py} & Prop.~\ref{prop:core} and its inputs & see below & $L\le5$, $s\le4$, $N\le7$\\
\texttt{t11\_sharp.py} & $r\ge f+2-\delta$; 2 controls & $160/160$, $368/368$; green & $\lambda\le5$, $f\le3$\\
\bottomrule
\end{tabular}
\end{center}

\texttt{t10\_proofinputs.py} is a positive control on the \emph{proof}, not only on the
statement. It forms the numerator $P=ZD$ in closed form using the exact cofactors
$(x_\alpha^2-x_\beta^2)/(\varepsilon_\alpha x_\alpha+\varepsilon_\beta x_\beta)
=\varepsilon_\alpha(x_\alpha-\varepsilon_\alpha\varepsilon_\beta x_\beta)$ and
$(c_p^2-x_\alpha^2)/(\varepsilon_\alpha x_\alpha+c_p)=c_p-\varepsilon_\alpha x_\alpha$,
and evaluates it exactly at random integer points. On the cells where $P\neq0$ --- so
that the checks are not vacuous --- it confirms Step 1's bound
$\deg_{x_1}P\le N-1$ (obtained by finite differences in $x_1$) and Step 2's conclusion
that $P$ vanishes on every hyperplane $x_\alpha=\pm x_\beta$, with \textbf{0 violations
of either} across all such cells. It also confirms that $P=0$ exactly on the $20$ cells
satisfying (i) or (ii) and nowhere else, and that $2(L-1)>L+s-1\iff L\ge s+2$ and
$2(L-1)+1>L+s-1\iff L\ge s+1$ with no mismatch over $L\le6$, $s\le5$.

\section{Grades}

\begin{center}
\small
\begin{tabular}{lll}
\toprule
statement & grade & basis\\
\midrule
Lemma~\ref{lem:parity}, Cor.~\ref{cor:dead} & \texttt{proved} & two-line proofs; $216/216$\\
Lemma~\ref{lem:bridge} (orbit factorisation) & \texttt{proved} & factor-by-factor; $63/63$ orbits\\
Lemma~\ref{lem:lin} (linearisation) & \texttt{proved} & one line\\
Lemma~\ref{lem:Zodd} (parity of $Z$) & \texttt{proved} & reindexing $\varepsilon_1\mapsto-\varepsilon_1$\\
Prop.~\ref{prop:core} (core identity, both cases) & \texttt{proved} & degree--divisibility squeeze\\
Thm.~\ref{thm:main} ($r\ge f+2-\delta$) & \texttt{proved} & Prop.~\ref{prop:core} + density\\
Cor.~\ref{cor:scope} & \texttt{proved} & specialisation of Thm.~\ref{thm:main}\\
Thm.~\ref{thm:refute} (converse false) & \texttt{proved} & 4 exact witnesses, two engines\\
Prop.~\ref{prop:global} ($f=0$ sign) & \texttt{proved} & reindexing; $140/140$\\
Cor.~\ref{cor:corrected}, second bullet & \texttt{computed} & $41$ cells $+$ $280$ control draws\\
Q578 parity profiles & \texttt{conjecture} & $16$ patterns at $r=4$ only\\
\bottomrule
\end{tabular}
\end{center}

No claim in this paper rests on any external source: every input is either proved here
or is one of my own earlier results (\texttt{thm:alt} 2026-10-22, the pair form of $B$
and \texttt{thm:r1} 2026-10-24 cycle~1), each re-verified by this session's independent
engine.

\end{document}
